\newpage
\setcounter{subsection}{0}
\section{CONTEXTO DEL PROYECTO}

\subsection{Trasfondo}

\subsubsection{Introducción a la Tecnología Blockchain y Contratos Inteligentes}

Una cadena de bloques (\textit{blockchain}) se define formalmente como una máquina de estados replicada y de transición discreta, cuyo estado global $\mathcal{S}$ evoluciona mediante la aplicación sucesiva de transiciones de estado transaccionales agrupadas en bloques ordenados cronológicamente. Matemáticamente, el estado del sistema en el instante discreto $t$ se formula como $\mathcal{S}_t = \Upsilon(\mathcal{S}_{t-1}, B_t)$, donde $\Upsilon$ representa la función de transición de estado global y $B_t$ es el bloque validado e incorporado en el registro histórico. Cada bloque $B_t$ constituye una estructura de datos inmutable enlazada criptográficamente con su predecesor inmediato a través de la inclusión del valor hash del encabezado anterior, denotado como $H(B_{t-1})$:
\begin{equation}
B_t = \left( H(B_{t-1}), \, \text{MerkleRoot}(T_t), \, \text{Timestamp}, \, \text{ConsensusMetadata} \right)
\end{equation}
donde $H: \{0,1\}^* \to \{0,1\}^{256}$ es una función hash criptográfica unidireccional y resistente a colisiones (como la familia SHA-2 o SHA-3), y $\text{MerkleRoot}(T_t)$ es la raíz del árbol de Merkle asociado al conjunto de transacciones $T_t = \{t_1, t_2, \dots, t_n\}$ incorporadas en dicho bloque. Esta disposición de punteros criptográficos garantiza la inmutabilidad y la auditabilidad eficiente del sistema: cualquier alteración retrospectiva en una transacción elemental $t_i$ se propaga recursivamente afectando la raíz de Merkle, los hashes intermedios y, en última instancia, el hash final del bloque correspondiente $B_t$, invalidando la consistencia de todos los bloques subsecuentes en la red distribuida.

Esta arquitectura de registro distribuido resuelve de manera determinística el problema clásico del doble gasto (\textit{double-spending}) sin requerir una entidad central de compensación. En el ámbito de los activos digitales tradicionales, la duplicación ilimitada y sin costo de archivos impide el establecimiento directo de escasez digital. El doble gasto se presenta cuando una entidad maliciosa intenta transferir de forma simultánea un mismo recurso financiero digital a dos destinatarios distintos. La tecnología blockchain mitiga este vector de ataque imponiendo un orden cronológico y global inalterable de las transacciones a través de un protocolo de consenso descentralizado. De este modo, solo la primera transacción que cumple con las condiciones de validez sobre el estado actual $\mathcal{S}_{t-1}$ es procesada por los nodos de la red, mientras que cualquier transacción competidora subsiguiente que intente gastar el mismo activo es descartada determinísticamente al violar las reglas de balance y propiedad del estado resultante.

Históricamente, el concepto de cadena de bloques y su primera implementación funcional se remontan a la publicación del documento blanco firmado bajo el seudónimo de Satoshi Nakamoto en el año 2008, titulado \textit{«Bitcoin: A Peer-to-Peer Electronic Cash System»} \cite{nakamoto2008bitcoin}. El surgimiento de esta tecnología responde a la necesidad estructural de prescindir de intermediarios centralizados de confianza (como bancos centrales, pasarelas de pago y cámaras de compensación) en las transacciones de valor en línea. Estos intermediarios tradicionales centralizan el control del flujo financiero, imponen estructuras tarifarias ineficientes y configuran puntos únicos de falla (*single points of failure*) vulnerables a la censura, ataques cibernéticos o inestabilidad sistémica.

Para ilustrar de forma clara e intuitiva el principio de inmutabilidad distribuida, se puede plantear una analogía técnica basada en un libro contable físico compartido. Imagínese una red de pares en la que no existe un escribano central único que custodie las cuentas. En su lugar, cada miembro de la red posee una copia idéntica y sincronizada del libro. Cada vez que se desea registrar un conjunto de transacciones (una página del libro), se debe aplicar una regla de validación matemática uniforme evaluada de forma independiente por todos los participantes. Si un actor malicioso altera de forma unilateral una línea en una página de su copia local, la discrepancia matemática en la firma acumulada y en el hash de enlace con las páginas posteriores se hace evidente de inmediato. Al contrastar copias, el resto de los participantes identificará la incongruencia y rechazará la versión fraudulenta del registro, manteniendo la integridad del estado global a través de la redundancia y el consenso distribuido.

Sobre esta base, los contratos inteligentes (*smart contracts*), propuestos conceptualmente por Nick Szabo en 1997 \cite{szabo1997formalizing} y popularizados con el despliegue de la red Ethereum en 2014 \cite{wood2014ethereum}, representan programas con estado persistente y lógica determinística autoejecutable. Estos contratos se despliegan directamente en el registro criptográfico y son evaluados por todos los validadores de la red como parte de la función de transición $\Upsilon$. Esto garantiza que sus reglas y compromisos se ejecuten exactamente según lo codificado, eliminando el riesgo de contraparte y la necesidad de supervisión manual externa.

\subsubsection{La Red Stellar}

Stellar es una red pública descentralizada orientada a la transferencia y movimiento eficiente de valor digital, a la tokenización de activos del mundo real (\textit{Real World Assets} o RWA) y a la inclusión financiera a escala global. Co-fundada en el año 2014 por Jed McCaleb y Joyce Kim bajo la gestión de la Stellar Development Foundation (SDF), la plataforma fue concebida con el propósito de mitigar las fricciones existentes en los sistemas de pagos transfronterizos tradicionales, los cuales se caracterizan por una alta latencia de liquidación, elevados costos de corresponsalía bancaria y una marcada falta de interoperabilidad.

A nivel de arquitectura de red y especificación operativa, Stellar destaca por las siguientes características técnicas clave:
\begin{itemize}\setlength{\itemsep}{2pt}
    \item \textbf{Eficiencia Energética y de Cómputo:} Al no depender de algoritmos de prueba de trabajo que exijan potencia de procesamiento masiva de carácter destructivo (como la minería ASIC de Bitcoin), Stellar opera mediante un protocolo de votación federada. El costo de procesamiento computacional por nodo es equivalente al de un servidor web estándar, reduciendo el consumo de energía y la huella de carbono del ledger global a valores insignificantes.
    \item \textbf{Baja Latencia de Cierre y Finalidad:} El tiempo de cierre de registro (\textit{ledger close time}) promedio de Stellar oscila entre los $3$ y $5$ segundos. Esta característica provee una finalidad transaccional absoluta e inmediata (*immediate finality*), elminando la necesidad de esperar confirmaciones heurísticas probabilísticas tras varias horas de computación, lo que resulta fundamental para procesos de emisión y validación de boletería digital en tiempo real.
    \item \textbf{Costos Operativos Insignificantes:} El costo base por operación en la red está configurado nominalmente en $100$ stroops (donde $1\text{ XLM} = 10^{7}\text{ stroops}$), lo cual equivale aproximadamente a $0{,}02\text{ COP}$ a la tasa de referencia histórica de conversión. Esta tarifa nominal actúa como un desincentivo criptoeconómico contra ataques de denegación de servicio por saturación de transacciones (DoS), permitiendo a la vez un volumen masivo de operaciones comerciales y de reventa sin penalizar económicamente a los usuarios del ecosistema.
\end{itemize}

\subsubsection{Protocolo de Consenso Stellar (SCP) y Acuerdo Bizantino Federado (FBA)}

El mecanismo de acuerdo y validación de la red Stellar está instrumentado bajo el Protocolo de Consenso Stellar (SCP, \textit{Stellar Consensus Protocol}) diseñado por David Mazières en 2015 \cite{mazieres2015stellar}. El SCP constituye la primera formulación académica y comercialmente implementada del paradigma del Acuerdo Bizantino Federado (FBA, \textit{Federated Byzantine Agreement}). 

Los sistemas clásicos de consenso tolerantes a fallas bizantinas (BFT), como el célebre protocolo PBFT (*Practical Byzantine Fault Tolerance*) \cite{castro1999practical}, imponen un modelo de membresía estática y cerrada. En estos esquemas, la participación requiere que una autoridad centralizada o un consorcio defina formalmente un conjunto estático de validadores $\mathbf{R}$. Todos los nodos de la red deben conocer y concordar unívocamente sobre la membresía del sistema en todo instante. Este enfoque limita drásticamente la descentralización al restringir el acceso libre de nodos e induce una penalización severa en el rendimiento debido a la complejidad de mensajería bidireccional exhaustiva del orden de $\mathcal{O}(|\mathbf{R}|^2)$ en su fase de confirmación.

En contraste, el FBA rompe con el supuesto de membresía cerrada y estática. En una red FBA, no existe una lista centralizada de validadores aprobados; cualquier nodo puede unirse al sistema y participar en el proceso de consenso de manera orgánica. Para alcanzar el acuerdo sin un listado global, el FBA introduce los siguientes conceptos matemáticos fundamentales:
\begin{itemize}\setlength{\itemsep}{2pt}
    \item \textbf{Rebanada de Quórum (\textit{Quorum Slice}):} Para un nodo individual $v \in V$, donde $V$ representa el conjunto de todos los nodos de la red, una rebanada de quórum $S(v)$ es un subconjunto no vacío de nodos tal que $v \in S(v) \subseteq V$, en el cual $v$ deposita su confianza de forma soberana e independiente para ratificar una transacción o transición de estado. Cada nodo puede definir de forma autónoma una o múltiples rebanadas de quórum alternativas (por ejemplo, basadas en organizaciones independientes o nodos altamente reputados) para estructurar su topología de confianza.
    \item \textbf{Quórum:} Un subconjunto de nodos $U \subseteq V$ constituye un quórum si es no vacío y para cada miembro del subconjunto $u \in U$, existe al menos una rebanada de quórum $S(u)$ tal que $S(u) \subseteq U$. De esta forma, las decisiones de aceptación de estado no se determinan por una mayoría absoluta y global, sino de manera inductiva a través de la intersección local y superpuesta de las rebanadas de quórum declaradas por cada participante.
\end{itemize}

Bajo el marco del Teorema CAP, que establece los límites fundamentales de consistencia (*Safety*), disponibilidad (*Liveness*) y tolerancia a la partición de red en sistemas distribuidos, el SCP prioriza strictly la consistencia sobre la disponibilidad. En la práctica, ante un fallo catastrófico de conectividad, un ataque bizantino masivo o una configuración errónea de confianza que impida la convergencia y la intersección de quórums, el protocolo prefiere congelar el avance de la red y suspender la confirmación de transacciones (*halt*) antes que permitir la generación de una bifurcación histórica en el registro. Esta priorización es la única coherente con sistemas de transferencia de activos financieros y boletería digital: la suspensión temporal del servicio resulta preferible a la confirmación de transacciones divergentes que posibiliten que una misma entrada digital sea asignada de manera definitiva a múltiples compradores concurrentes.

A nivel energético y de gobernanza, el SCP exhibe una eficiencia superior en comparación con los esquemas de consenso preponderantes en la industria:
\begin{itemize}\setlength{\itemsep}{2pt}
    \item \textbf{Proof of Work (PoW):} El consenso se alcanza mediante una carrera computacional competitiva de fuerza bruta para encontrar el hash de un bloque que cumpla con una dificultad dinámica. Esto implica un desperdicio energético masivo, del orden de teravatios-hora a nivel global, y restringe el rendimiento a unas pocas transacciones por segundo debido a los tiempos de dificultad adaptativa.
    \item \textbf{Proof of Stake (PoS):} Si bien mitiga el consumo de energía, la votación está directamente correlacionada con el capital criptográfico depositado o bloqueado (*staking*). Este modelo introduce riesgos de monopolio y consolidación plutocrática del control de la red por parte de los tenedores mayoritarios de tokens, además de requerir algoritmos complejos para resolver escenarios de bifurcación sin costo (*nothing-at-stake*) y mecanismos de penalización monetaria extrema (*slashing*).
    \item \textbf{Stellar Consensus Protocol (SCP):} El acuerdo se construye puramente sobre el intercambio de mensajes criptográficamente firmados entre nodos validados según la topología descentralizada de las rebanadas de confianza. Al no requerir computación intensiva de hashing ni bloqueos financieros coercitivos de capital para la ponderación del voto, el consumo energético de la red Stellar se asemeja al de la infraestructura de mensajería de un sistema web tradicional de alta eficiencia, asegurando rapidez y un costo de carbono mínimo.
\end{itemize}

\subsubsection{Tokens Estables (Stablecoins) y Bases Matemáticas del Consenso}

Las monedas estables (\textit{stablecoins}) son representaciones digitales de valor emitidas sobre un registro blockchain cuyo valor de intercambio comercial se encuentra referenciado o respaldado en paridad unitaria ($1:1$) con respecto a un activo tradicional de baja volatilidad, comúnmente una moneda fiduciaria de curso legal como el dólar estadounidense (USD). En la arquitectura Web2.5 planteada en este trabajo, el empleo de tokens estables como USDC (desplegado nativamente en Stellar mediante el \textit{Stellar Asset Contract} o SAC) resulta un pilar conceptual fundamental para solventar la alta variabilidad del activo nativo de la red (XLM). La volatilidad inherente de los criptoactivos nativos impediría la adopción del sistema por parte de operadores logísticos y usuarios finales, debido al riesgo cambiario implícito en los intervalos temporales de reventa. El uso de un token digital convertible a paridad fiduciaria asegura la certidumbre económica del precio final y de las comisiones del mercado secundario.

La estabilidad y consistencia global del Acuerdo Bizantino Federado descansa sobre la base matemática de la \textbf{intersección de quórums}. Decimos formalmente que un sistema de quórums $\mathcal{Q}$ para una red federada $V$ posee la propiedad de intersección de quórums si para todo par de quórums $U_1, U_2 \in \mathcal{Q}$, se cumple que:
\begin{equation}
U_1 \cap U_2 \neq \emptyset
\end{equation}
y dicha intersección contiene al menos un nodo correcto (no defectuoso ni bizantino). Esta propiedad garantiza matemáticamente que los nodos honestos de la red no podrán validar transacciones contradictorias en registros paralelos, dado que cualquier decisión final de consenso en un quórum requiere la ratificación de nodos compartidos con el otro quórum. La conformación adecuada de las rebanadas de quórum por parte de los operadores de nodos asegura que la red mantenga la intersección de quórums de forma dinámica, tolerando hasta un tercio ($f < n/3$) de fallos bizantinos locales en las subestructuras de confianza sin comprometer la consistencia global del ledger distribuido.


\subsubsection{Soroban, NFTs y tokens de pago}

Soroban es la plataforma de contratos inteligentes de Stellar. Los contratos se escriben en Rust contra el crate \texttt{soroban-sdk} (versión 23 en este trabajo) y se compilan al target \texttt{wasm32v1-none}. El modelo de ejecución es determinístico, sin acceso a fuentes no reproducibles (relojes locales, aleatoriedad no derivada del ledger, llamadas al sistema operativo).

Soroban gestiona tres tipos de almacenamiento: \texttt{Persistent}, \texttt{Temporary} e \texttt{Instance}, cada uno con un costo de \textit{rent} distinto. La emisión de eventos se realiza mediante \texttt{events().publish()}: cada evento contiene \textit{topics} indexables y un \textit{body} serializable, y se almacena en el meta del ledger, lo que permite a los indexadores externos consumirlo sin afectar el costo de \textit{rent} del contrato.

El estándar NFT en Soroban define una representación de activos no fungibles con identificador único, propietario tipo \texttt{Address} y reglas de transferencia controladas por el contrato emisor. En este trabajo se extiende este estándar con un esquema de versionado: cada boleto se identifica por la tupla (\texttt{ticket\_root\_id}, \texttt{version}), donde \texttt{version} se incrementa en cada operación de reventa mediante un patrón \textit{burn/remint}.

El activo nativo de la red, XLM, y cualquier activo emitido en Stellar se exponen como tokens compatibles con la interfaz estándar de Soroban a través del \textit{Stellar Asset Contract} (SAC). Esto permite que las operaciones de pago de la compra de reventa se ejecuten dentro del mismo flujo de invocación del contrato del evento, sin transacciones separadas para el pago y la transferencia de propiedad.

\subsection{Análisis del Contexto}

El análisis comparativo de soluciones del mercado y de la literatura técnica permite identificar tres categorías de antecedentes: (i) plataformas tradicionales de reventa basadas en QR estáticos y conciliación bancaria; (ii) soluciones NFT desplegadas sobre Ethereum 1.0 o sus capas 2; y (iii) propuestas académicas de \textit{ticketing} distribuido sin liquidación atómica. La Tabla~\ref{tab:analisis-contexto} contrasta los principales parámetros técnicos de cada categoría contra la propuesta de este trabajo.

\begin{table}[H]
\centering
\small
\caption{Análisis comparativo de soluciones del contexto.}
\label{tab:analisis-contexto}
\begin{tabular}{p{3.2cm}p{2.7cm}p{2.7cm}p{2.7cm}p{2.7cm}}
\toprule
\textbf{Parámetro} & \textbf{Tradicional} & \textbf{NFT en ETH 1.0} & \textbf{Académica} & \textbf{Propuesta} \\
\midrule
Finalidad & 72--120 h & 1--6 min & Variable & 3--5 s \\
Costo por op. & 1{,}5--3{,}5\,\% + fija & USD 5--50 (gas) & N/A & $\approx$\,0{,}02 COP \\
Capacidad & --- & 15--30 TPS & --- & 1\,000--3\,000 TPS \\
Atomicidad pago + propiedad & No & Parcial & No & Sí \\
Trazabilidad pública & No & Sí & Parcial & Sí \\
Unicidad por construcción & No & Sí & Variable & Sí (\textit{burn/remint}) \\
\bottomrule
\end{tabular}
\end{table}

La principal ventaja de la propuesta frente al esquema NFT clásico sobre Ethereum es la liquidación atómica de pago, comisiones y propiedad dentro de una sola transacción, sin requerir patrones de \textit{commit-reveal} o \textit{escrow} externos. Frente a las plataformas tradicionales, la unicidad del NFT con versionado y el patrón \textit{burn/remint} eliminan por construcción la coexistencia de dos copias válidas del mismo activo, lo cual no es alcanzable mediante validación posterior de QR estáticos.

Como referencia normativa local, la Ley 1493 de 2011 regula los espectáculos públicos de las artes escénicas en Colombia y establece una contribución parafiscal del 10\,\% sobre la boletería cuyo precio individual sea igual o superior a tres UVT. El modelo propuesto registra el precio efectivo on-chain en cada operación, lo que habilita el cálculo determinístico de la contribución mediante una agregación sobre el historial indexado. Este mapeo se utiliza como referencia académica y no constituye declaración de cumplimiento regulatorio productivo.
